\documentclass[11pt]{article}
\usepackage{fontspec,amsmath,amssymb,amsthm,hyperref}
\usepackage[margin=1in]{geometry}
\setlength{\emergencystretch}{3em}
\title{Proposition 1: scalar warmup, variable steps and iteration count}
\author{LangevinLean mathematical exposition}
\date{}
\begin{document}\maketitle
\noindent Source module: \texttt{AverageLMC/Complexity.lean}. Compile this proof document with XeLaTeX or LuaLaTeX. The accompanying HTML pairs the mathematical argument with the exact Lean statements and records the scope of the formalization.


Use \(a=\kappa_{\mathrm{av}}+1>0\), \(\kappa=\kappa_{\max}\ge1\), and \(\varepsilon_0>0\). The paper's exact positive logarithm, warmup ceiling, decreasing-step schedule and accuracy ceiling are represented explicitly. Given \(0\le r_0\le\varepsilon_0^2\) and the normalized scalar one-step recurrence, the module proves \(r_k\le a/(\kappa+k-k_0)\) for \(k\ge k_0\), then \(r_k\le\varepsilon^2\) at the stated iteration count. All stochastic, coupling and Wasserstein facts needed to establish that recurrence remain external proof obligations.


\section{Encode the exact positive logarithm and two-phase schedule}


The definitions are \[\log_+(z)=\max\{\log z,0\},\qquad k_0=\left\lceil\frac\kappa2\log_+\!\left(\frac{2\kappa\varepsilon_0^2}{a}\right)\right\rceil,\qquad t_k=\frac1{\kappa+(k-k_0)_+}.\] In Lean, the ceiling is a natural-number ceiling and \((k-k_0)_+\) is represented by subtraction in the natural numbers, cast afterward to \(\mathbb R\). Thus the difference is exactly zero before warmup ends; for \(k\ge k_0\) it is the ordinary integer difference. The dimensionless step \(t_k=mh_k\) corresponds to the paper's \(h_k=1/(M_{\max}+m(k-k_0)_+)\). These are definitions, with no hidden decay or convergence assumptions.



\noindent Lean declarations: \texttt{positiveLog}, \texttt{warmup}, \texttt{variableStep}.

\section{Verify the smoothness restriction for the full schedule}


For \(\kappa>0\), the cast natural difference is nonnegative, so \(\kappa+(k-k_0)_+\ge\kappa>0\). Positivity of the reciprocal and monotonicity of reciprocals give \[0<t_k\le1/\kappa.\] When \(\kappa=\kappa_{\max}=M_{\max}/m\), this is the arithmetic behind \(h_kM_{\max}\le1\). It verifies the schedule's restriction for every step, including those in the warmup phase; it does not assert the one-step stochastic inequality itself.



\noindent Lean declarations: \texttt{variableStep\_le}.

\section{Use the warmup ceiling to reduce the initial error}


Let \(R=2\kappa\varepsilon_0^2/a>0\). The natural ceiling gives \[k_0\ge\frac\kappa2\max\{\log R,0\}\ge\frac\kappa2\log R.\] Multiplying and dividing by the positive \(\kappa\) yields \(-2k_0/\kappa\le-\log R\). Exponentiate this inequality and multiply by the nonnegative \(\varepsilon_0^2\): \[e^{-2k_0/\kappa}\varepsilon_0^2\le e^{-\log R}\varepsilon_0^2=\frac a{2\kappa}.\] The final equality uses \(e^{\log R}=R\) on its verified positive domain and cancels nonzero parameters. The same proof covers \(R\le1\), where the positive logarithm is zero and the warmup length can vanish.



\noindent Lean declarations: \texttt{warmup\_decay}.

\section{Join warmup to the harmonic tail}


Assume \(0\le r_0\le\varepsilon_0^2\) and, for every natural \(k\), \[r_{k+1}\le\frac{1-t_k}{1+t_k}r_k+\frac{at_k^2}{1+t_k},\qquad t_k=\frac1{\kappa+(k-k_0)_+}.\] Before \(k_0\), natural subtraction gives \(t_k=1/\kappa\). The finite-horizon constant-step theorem therefore gives \[r_{k_0}\le e^{-2k_0/\kappa}r_0+\frac a{2\kappa}\le\frac a\kappa,\] where the second inequality uses the initialization bound and the warmup-decay lemma. Shift the sequence to \(s_n=r_{k_0+n}\). Cancellation of the natural subtraction gives \(t_{k_0+n}=1/(\kappa+n)\), so the harmonic tail theorem applies and proves \(s_n\le a/(\kappa+n)\). For any \(k\ge k_0\), choose \(n=k-k_0\) and use \(k_0+(k-k_0)=k\), obtaining \[r_k\le\frac a{\kappa+(k-k_0)}.\] This is precisely the scalar calculation in the proof of Proposition 1, with the recurrence kept as an explicit hypothesis rather than disguised as a convexity or smoothness assumption.



\noindent Lean declarations: \texttt{variable\_step\_normalized}.

\section{Verify the accuracy ceiling and conclude the conditional rate}


The exact iteration threshold is \[k_\varepsilon=k_0+\left\lceil\max\left\{\frac a{\varepsilon^2}-\kappa,0\right\}\right\rceil.\] Suppose \(\varepsilon>0\) and \(k\ge k_\varepsilon\). Natural arithmetic first gives \(k\ge k_0\) and then bounds \(k-k_0\) below by that ceiling. After casting to \(\mathbb R\), \[k-k_0\ge\max\{a/\varepsilon^2-\kappa,0\}\ge a/\varepsilon^2-\kappa.\] Thus \(\kappa+(k-k_0)\ge a/\varepsilon^2\). Both denominators used in rearrangement are positive, proving \[\frac a{\kappa+(k-k_0)}\le\varepsilon^2.\] The arithmetic envelope theorem proves this implication without a recurrence. The final accuracy theorem combines it with the preceding conditional scalar Proposition 1 to obtain \(r_k\le\varepsilon^2\). It proves no independent assertion that actual LMC errors obey the required recurrence.



\noindent Lean declarations: \texttt{variableIterations}, \texttt{variable\_iterations\_envelope}, \texttt{variable\_step\_accuracy}.

\end{document}
